% ENEM 2019-2025 — 6 questões MCQ — Aritmética II
% Fonte: lista "Aritmética II" do site projetoagathaedu.com.br
% Omitidas: Q1,Q2,Q4,Q6,Q7,Q8,Q12,Q13,Q14,Q15,Q16,Q17 (imagens ou dados faltando)

---
tipo: Múltipla Escolha
dificuldade: Média
nivel: médio
disciplina: Matemática
assunto: Progressão Geométrica
gabarito: C
tags: [ENEM, PG, progressão geométrica, tempo, sequência]
fonte: concurso
concurso: ENEM
banca: INEP
ano: 2019
numero: 3
---
\question
Alguns modelos de rádios automotivos estão protegidos por um código de segurança. Para ativar o sistema de áudio, deve-se digitar o código secreto composto por quatro algarismos. No primeiro caso de erro na digitação, a pessoa deve esperar 60 segundos para digitar o código novamente. O tempo de espera duplica, em relação ao tempo de espera anterior, a cada digitação errada. Uma pessoa conseguiu ativar o rádio somente na quarta tentativa, sendo de 30 segundos o tempo gasto para digitação do código secreto a cada tentativa. Nos casos da digitação incorreta, ela iniciou a nova tentativa imediatamente após a liberação do sistema de espera.

O tempo total, em segundo, gasto por essa pessoa para ativar o rádio foi igual a
\begin{choices}
\choice 300.
\choice 420.
\CorrectChoice 540.
\choice 660.
\choice 1.020.
\end{choices}

---
tipo: Múltipla Escolha
dificuldade: Média
nivel: médio
disciplina: Matemática
assunto: Números Inteiros
gabarito: C
tags: [ENEM, números inteiros, operações, elevador, andares]
fonte: concurso
concurso: ENEM
banca: INEP
ano: 2025
numero: 5
---
\question
Um edifício tem a numeração dos andares iniciando no térreo (T), e continuando com primeiro, segundo, terceiro, ..., até o último andar. Uma criança entrou no elevador e, tocando no painel, seguiu uma sequência de andares, parando, abrindo e fechando a porta em diversos andares. A partir de onde entrou a criança, o elevador subiu sete andares, em seguida desceu dez, desceu mais treze, subiu nove, desceu quatro e parou no quinto andar, finalizando a sequência. Considere que, no trajeto seguido pela criança, o elevador parou uma vez no último andar do edifício.

De acordo com as informações dadas, o último andar do edifício é o
\begin{choices}
\choice 16º
\choice 22º
\CorrectChoice 23º
\choice 25º
\choice 32º
\end{choices}

---
tipo: Múltipla Escolha
dificuldade: Média
nivel: médio
disciplina: Matemática
assunto: Progressão Aritmética
gabarito: C
tags: [ENEM, PA, progressão aritmética, postes, iluminação]
fonte: concurso
concurso: ENEM
banca: INEP
ano: 2025
numero: 9
---
\question
A prefeitura de um pequeno município do interior decide colocar postes para iluminação ao longo de uma estrada retilínea, que inicia em uma praça central e termina numa fazenda na zona rural. Como a praça já possui iluminação, o primeiro poste será colocado a 80 metros da praça, o segundo, a 100 metros, o terceiro, a 120 metros, e assim sucessivamente, mantendo-se sempre uma distância de vinte metros entre os postes, até que o último poste seja colocado a uma distância de 1.380 metros da praça.

Se a prefeitura pode pagar, no máximo, R$ 8.000,00 por poste colocado, o maior valor que poderá gastar com a colocação desses postes é
\begin{choices}
\choice R$ 512.000,00
\choice R$ 520.000,00
\CorrectChoice R$ 528.000,00
\choice R$ 552.000,00
\choice R$ 584.000,00
\end{choices}

---
tipo: Múltipla Escolha
dificuldade: Média
nivel: médio
disciplina: Matemática
assunto: Progressão Geométrica
gabarito: E
tags: [ENEM, PG, progressão geométrica, torneio, eliminatória, tênis]
fonte: concurso
concurso: ENEM
banca: INEP
ano: 2025
numero: 10
---
\question
Torneios de tênis, em geral, são disputados em sistema de eliminatória simples. Nesse sistema, são disputadas partidas entre dois competidores, com a eliminação do perdedor e promoção do vencedor para a fase seguinte. Dessa forma, se na 1ª fase o torneio conta com $2n$ competidores, então na 2ª fase restarão $n$ competidores, e assim sucessivamente até a partida final.

Em um torneio de tênis, disputado nesse sistema, participam 128 tenistas.

Para se definir o campeão desse torneio, o número de partidas necessárias é dado por
\begin{choices}
\choice $2 \times 128$
\choice $64 + 32 + 16 + 8 + 4 + 2$
\choice $128 + 64 + 32 + 16 + 8 + 4 + 2 + 1$
\choice $128 + 64 + 32 + 16 + 8 + 4 + 2$
\CorrectChoice $64 + 32 + 16 + 8 + 4 + 2 + 1$
\end{choices}

---
tipo: Múltipla Escolha
dificuldade: Fácil
nivel: médio
disciplina: Matemática
assunto: Aritmética
gabarito: E
tags: [ENEM, aritmética, código, sistema de codificação, memorando]
fonte: concurso
concurso: ENEM
banca: INEP
ano: 2025
numero: 11
---
\question
Uma repartição pública possui um sistema que armazena em seu banco de dados todos os ofícios, memorandos e cartas enviados ao longo dos anos. Para organizar todo esse material e facilitar a localização no sistema, o computador utilizado pela repartição gera um código para cada documento, de forma que os oito primeiros dígitos indicam a data em que o documento foi emitido (DDMMAAAA), os dois dígitos seguintes indicam o tipo de documento (ofício: 01, memorando: 02 e carta: 03) e os três últimos dígitos indicam a ordem do documento. Por exemplo, o código 0703201201003 indica um ofício emitido no dia 7 de março de 2012, cuja ordem é 003. No dia 27 de janeiro de 2001, essa repartição pública emitiu o memorando de ordem 012 e o enviou aos seus funcionários.

O código gerado para esse memorando foi
\begin{choices}
\choice 0122701200102.
\choice 0201227012001.
\choice 0227012001012.
\choice 2701200101202.
\CorrectChoice 2701200102012.
\end{choices}

---
tipo: Múltipla Escolha
dificuldade: Média
nivel: médio
disciplina: Matemática
assunto: Aritmética
gabarito: C
tags: [ENEM, aritmética, custo, otimização, mirante, elevador, teleférico]
fonte: concurso
concurso: ENEM
banca: INEP
ano: 2025
numero: 18
---
\question
Em um parque há dois mirantes de alturas distintas que são acessados por elevador panorâmico. O topo do mirante 1 é acessado pelo elevador 1, enquanto que o topo do mirante 2 é acessado pelo elevador 2. Eles encontram-se a uma distância possível de ser percorrida a pé, e entre os mirantes há um teleférico que os liga, que pode ou não ser utilizado pelo visitante.

O acesso aos elevadores tem os seguintes custos:
\begin{itemize}
\item Subir pelo elevador 1: R\$ 0,15;
\item Subir pelo elevador 2: R\$ 1,80;
\item Descer pelo elevador 1: R\$ 0,10;
\item Descer pelo elevador 2: R\$ 2,30.
\end{itemize}

O custo da passagem do teleférico partindo do topo do mirante 1 para o topo do mirante 2 é de R\$ 2,00, e do topo do mirante 2 para o topo do mirante 1 é de R\$ 2,50.

Qual é o menor custo, em real, para uma pessoa visitar os topos dos dois mirantes e retornar ao solo?
\begin{choices}
\choice 2,25
\choice 3,90
\CorrectChoice 4,35
\choice 4,40
\choice 4,45
\end{choices}
